% Zone „Das kennst du schon“ – aus bank/scharen-von-geraden-und-ebenen/zone.jsonl (zusammenbau v0.1)
\einheitenkopf[zone][Kennst du schon]{Das kennst du schon}
\setcounter{aufgabe}{0}

%% TODO Titel: Ich-kann-Satz fehlt in der Bank (Platzhalter: Kettenname) – Nr. 1
%% TODO Anweisung: ein Satz über der ersten Teilaufgabe fehlt in der Bank – Nr. 1
\begin{aufgabe}{Punktprobe an Geraden- und Ebenengleichungen (einsetzen, Gleichung prüfen)}
\begin{teile}
\teil Liegt $P(1 | 2 | 3)$ in der Ebene $E: x_1 + x_2 + x_3 = 6$? \janein
\teil Liegt $P(5 | -1 | 4)$ auf der Geraden $g: \vec{x} = (1 | 1 | 2) + r \cdot (2 | -1 | 1)$? \janein
\end{teile}
\end{aufgabe}

%% TODO Titel: Ich-kann-Satz fehlt in der Bank (Platzhalter: Kettenname) – Nr. 2
%% TODO Anweisung: ein Satz über der ersten Teilaufgabe fehlt in der Bank – Nr. 2
\begin{aufgabe}{Normalenvektor ablesen und die Paarregeln anwenden (parallel: Skalarprodukt null; senkrecht: kollinear)}
\begin{teile}
\teil Normalenvektor von $E: 2x_1 - 3x_2 + x_3 = 5$? $\vec{n} = ($ \leerfeld | \leerfeld | \leerfeld $)$
\teil Ist eine Gerade mit Richtungsvektor $\vec{u} = (3 | 1 | 1)$ parallel zur Ebene $E: x_1 - 2x_2 - x_3 = 4$? \janein
\end{teile}
\end{aufgabe}

%% TODO Titel: Ich-kann-Satz fehlt in der Bank (Platzhalter: Kettenname) – Nr. 3
%% TODO Anweisung: ein Satz über der ersten Teilaufgabe fehlt in der Bank – Nr. 3
\begin{aufgabe}{Kollinearität prüfen und ein Gleichungssystem auf Widerspruch untersuchen}
\begin{teile}
\teil Sind $(1 | -2 | 3)$ und $(3 | -6 | 9)$ kollinear? \janein
\end{teile}
\begin{gleichungsraster}[2]
\gl{\text{Hat das System $r + s = 3$, $r - s = 1$, $2r + s = 4$ eine Lösung?}} & \\
\end{gleichungsraster}
\end{aufgabe}

%% TODO Titel: Ich-kann-Satz fehlt in der Bank (Platzhalter: Kettenname) – Nr. 4
%% TODO Anweisung: ein Satz über der ersten Teilaufgabe fehlt in der Bank – Nr. 4
\begin{aufgabe}{Lineare und quadratische Gleichungen sowie Ungleichungen im Parameter lösen (quadrieren, Beträge, Vorzeichenfälle)}
\begin{gleichungsraster}[2]
\gl{\text{Löse $3k - 2 = k + 6$.}} &
\gl{\text{Löse $k^2 - 2k = 8$.}} \\
\end{gleichungsraster}
\end{aufgabe}

%% TODO Titel: Ich-kann-Satz fehlt in der Bank (Platzhalter: Kettenname) – Nr. 5
%% TODO Anweisung: ein Satz über der ersten Teilaufgabe fehlt in der Bank – Nr. 5
\begin{aufgabe}{Winkelformeln (Kosinus für Vektorpaare und Ebenen, Sinus für Gerade–Ebene) und die Hessesche Normalform anwenden}
\begin{teile}
\teil Winkel zwischen $\vec{u} = (1 | 1 | 0)$ und $\vec{v} = (0 | 1 | 1)$? $\alpha = $ \leerfeld
\teil Winkel zwischen einer Geraden mit $\vec{u} = (0 | 3 | 3)$ und der Ebene $E: x_3 = 4$? $\alpha = $ \leerfeld
\end{teile}
\end{aufgabe}

%% TODO Titel: Ich-kann-Satz fehlt in der Bank (Platzhalter: Kettenname) – Nr. 6
%% TODO Anweisung: ein Satz über der ersten Teilaufgabe fehlt in der Bank – Nr. 6
\begin{aufgabe}{Schnittpunkte einer Ebene mit Kanten und Achsen berechnen und Schrägbilder lesen}
\begin{teile}
\teil Wo schneidet $E: x_1 + 2x_2 + 4x_3 = 8$ die $x_1$-Achse? $S($ \leerfeld | \leerfeld | \leerfeld $)$
\teil Wo schneidet $E: x_1 + x_2 + x_3 = 7$ die Kante von $P(2 | 2 | 0)$ bis $Q(2 | 2 | 4)$? $S($ \leerfeld | \leerfeld | \leerfeld $)$
\end{teile}
\end{aufgabe}

%% TODO Titel: Ich-kann-Satz fehlt in der Bank (Platzhalter: Kettenname) – Nr. 7
%% TODO Anweisung: ein Satz über der ersten Teilaufgabe fehlt in der Bank – Nr. 7
\begin{aufgabe}{Normalenvektor ablesen und die Paarregeln anwenden (parallel: Skalarprodukt null; senkrecht: kollinear) – Fehler finden}
\begin{teile}
\teil Finde den Fehler: Steht $g$ mit $\vec{u} = (1 | 1 | -2)$ senkrecht auf $E: 3x_1 - x_2 + x_3 = 0$? \rechnung{\vec{n} \cdot \vec{u} &= 3 - 1 - 2 = 0 && \Rightarrow g \text{ steht senkrecht auf } E}
\end{teile}
\end{aufgabe}

%% TODO Titel: Ich-kann-Satz fehlt in der Bank (Platzhalter: Kettenname) – Nr. 8
%% TODO Anweisung: ein Satz über der ersten Teilaufgabe fehlt in der Bank – Nr. 8
\begin{aufgabe}{Normalenvektor ablesen und die Paarregeln anwenden (parallel: Skalarprodukt null; senkrecht: kollinear) – selbst rechnen}
\begin{teile}
\teil Steht $g$ mit $\vec{u} = (-2 | 1 | 3)$ senkrecht auf $E: 4x_1 - 2x_2 - 6x_3 = 3$? \janein
\end{teile}
\end{aufgabe}
